\newpage
\section{Giełdy zdecentralizowane i~arbitraż --- przegląd zagadnień}

\subsection{Ethereum: bloki, transakcje, zdarzenia i~opłaty za gaz}

Ethereum jest zdecentralizowaną maszyną stanów. Stan globalny zmieniają wyłącznie transakcje
zgrupowane w~bloki, a~semantykę ich wykonania definiuje maszyna wirtualna EVM
(ang. \emph{Ethereum Virtual Machine})~\cite{ethereum_yellowpaper}. Kontrakt może podczas wykonania
wyemitować \emph{zdarzenie} (ang. \emph{event}), które trafia do paragonu transakcji jako
\emph{log}: adres kontraktu, do czterech indeksowanych tematów (ang. \emph{topics}) i~dane.
Węzły udostępniają metodę JSON-RPC \texttt{eth\_getLogs}, pobierającą logi z~zakresu bloków według
adresu i~tematu. Kontrakty giełd emitują zdarzenie po każdej zmianie rezerw, więc logi są
kompletnym i~weryfikowalnym zapisem historii puli. Na tej własności opiera się warstwa
pozyskiwania danych z~rozdziału~3.4.

Każda operacja EVM ma koszt wyrażony w~jednostkach \emph{gazu}. Nadawca deklaruje limit gazu
i~cenę za jednostkę. W~pierwotnym modelu opłat cena gazu była pojedynczą wartością
\texttt{gasPrice} ustalaną przez nadawcę; górnicy porządkowali transakcje malejąco według tej
ceny, co sprowadzało dostęp do bloku do jawnej licytacji. Zmiana EIP-1559~\cite{eip1559},
aktywowana w~bloku 12\,965\,000 (5~sierpnia 2021~r.), wprowadziła opłatę bazową
\texttt{baseFeePerGas}, ustalaną przez protokół i~spalaną, oraz napiwek
\texttt{maxPriorityFeePerGas} dla producenta bloku. Z~analizy Liu i~in.~\cite{liu_eip1559} wynika,
że EIP-1559 zmniejszyło rozrzut cen gazu wewnątrz bloku i~ułatwiło ich szacowanie, ale samego
poziomu opłat istotnie nie obniżyło. Dla tej pracy oznacza to dwie rzeczy: koszt transakcji
arbitrażowej trzeba szacować inaczej w~oknach sprzed i~po EIP-1559 (rozdział~3.4), a~cena gazu
pozostaje wielkością zmienną w~czasie i~jednym z~głównych składników kosztu arbitrażu.

\subsection{Automatyczni animatorzy rynku o~stałym iloczynie}

Giełda typu AMM zastępuje księgę zleceń pulą płynności. Kontrakt przechowuje rezerwy dwóch
tokenów i~wycenia wymianę deterministyczną funkcją tych rezerw; kontrahentem każdej transakcji
jest sama pula~\cite{xu_sok_amm}. W~protokole Uniswap~V2~\cite{uniswap_v2} funkcją tą jest stały
iloczyn. Dla rezerw $x$ i~$y$ tokenów $X$ i~$Y$ niezmiennik ma postać
\begin{equation}
  x \cdot y = k .
  \label{eq:cpmm}
\end{equation}
Cena chwilowa tokena $X$ wyrażona w~tokenie $Y$ wynika ze stosunku rezerw:
\begin{equation}
  p = \frac{y}{x}.
  \label{eq:spot}
\end{equation}
Wymiana $\Delta x$ jednostek $X$ przy prowizji $\varphi = 0{,}3\,\%$ zwraca
\begin{equation}
  \Delta y = \frac{(1-\varphi)\,\Delta x \cdot y}{x + (1-\varphi)\,\Delta x}
           = \frac{997\,\Delta x\, y}{1000\,x + 997\,\Delta x},
  \label{eq:amountout}
\end{equation}
zgodnie z~funkcją \texttt{getAmountOut} w~kontraktach Uniswap~V2. Z~równań \eqref{eq:cpmm} i~\eqref{eq:amountout} widać, że cena efektywna rośnie z~wielkością transakcji (poślizg cenowy),
a~prowizja zostaje w~puli i~powiększa $k$. Po każdej zmianie stanu kontrakt pary emituje
zdarzenie \texttt{Sync} z~nowymi rezerwami, a~przy wymianie dodatkowo \texttt{Swap} z~kwotami.
Cenę~\eqref{eq:spot} w~dowolnym bloku można więc odtworzyć z~samych logów.

Sushiswap powstał w~2020~r. jako kopia (ang. \emph{fork}) kontraktów Uniswap~V2 z~własnymi pulami
płynności. Oba protokoły mają identyczną formułę~\eqref{eq:amountout}, te same sygnatury zdarzeń
i~tę samą prowizję, więc porównanie ich cen sprowadza się do porównania rezerw w~tym samym
bloku. Systematyczny przegląd protokołów AMM, także tych o~stałej średniej i~o~skoncentrowanej
płynności, zawiera praca Xu i~in.~\cite{xu_sok_amm}.

Angeris i~in.~\cite{angeris_uniswap} pokazali, że w~rynku o~stałym iloczynie każde odchylenie
ceny~\eqref{eq:spot} od ceny referencyjnej tworzy zyskowną okazję dla arbitrażysty, a~zysk ten
maksymalizuje transakcja zamykająca odchylenie do poziomu prowizji. Przy dostatecznej
aktywności arbitrażystów cena puli musi zatem śledzić rynek referencyjny. Obserwowana
rozbieżność między dwiema pulami świadczy o~czymś przeciwnym: w~danym bloku arbitraż nie został
jeszcze wykonany.

\subsection{Arbitraż między pulami i~czynniki jego wykonalności}

Niech dwie pule tej samej pary tokenów, $A$ i~$B$, mają ceny chwilowe $p_A$ i~$p_B$. Względną
rozbieżność cenową definiuje się jako
\begin{equation}
  S = \frac{|p_A - p_B|}{\min(p_A, p_B)} \cdot 100\,\%.
  \label{eq:spread}
\end{equation}
Arbitraż dwupulowy to zakup tokena w~puli tańszej i~sprzedaż w~droższej w~jednej, atomowej
transakcji. Każda z~dwóch wymian jest obciążona prowizją $\varphi$, więc zysk brutto pojawia się
dopiero, gdy rozbieżność przekracza koszt prowizji obu nóg, $(1-\varphi)^{-2} - 1 \approx
0{,}60\,\%$. W~pracy przyjęto próg $0{,}65\,\%$ (rozdział~3.5). Dla ustalonego stanu obu puli
istnieje wielkość transakcji maksymalizująca zysk brutto; wynika ona z~równania
\eqref{eq:amountout} i~jest wyprowadzona w~rozdziale~3.5.

Rozbieżność powyżej progu nie przesądza jednak o~wykonalności. Literatura wskazuje trzy grupy
czynników. Pierwszą jest koszt wykonania: transakcja arbitrażowa zużywa gaz na dwie wymiany
i~transfery, a~przy wysokiej cenie gazu ten koszt stały może przewyższyć zysk brutto. Drugą jest
płynność. Zysk z~rozbieżności ogranicza poślizg, więc w~płytkich pulach optymalna transakcja
jest mała, a~zysk absolutny niewielki~\cite{angeris_uniswap}. Trzecią grupą jest konkurencja
i~kolejność transakcji: okazję widzą wszyscy obserwatorzy sieci, a~o~tym, kto ją zrealizuje,
decyduje kolejność w~bloku.

Ostatni czynnik jest przedmiotem badań nad wartością możliwą do wydobycia przez producenta
bloku (ang. \emph{miner/maximal extractable value}, MEV). Daian i~in.~\cite{daian_flashboys}
udokumentowali boty arbitrażowe konkurujące o~pierwszeństwo w~bloku przez licytację ceny gazu
(ang. \emph{priority gas auctions}); znaczna część zysku z~arbitrażu trafia wtedy do górników,
a~sama konkurencja destabilizuje konsensus. Zhou i~in.~\cite{zhou_hft_dex} sformalizowali ataki na
kolejność transakcji na giełdach AMM (\emph{front-running}, \emph{sandwich}), a~Qin
i~in.~\cite{qin_bev} oszacowali łączną skalę wartości wydobywanej z~DeFi, w~tym z~arbitrażu.
Wniosek istotny dla tej pracy: aktywność botów i~podwyższona cena gazu w~danym momencie
zmniejszają szansę, że obserwator rozbieżności skonsumuje ją jako pierwszy. Te dwie wielkości
tworzą cechę ryzyka~$M$ w~rozdziale~3.5. Odkąd Flashbots udostępnili prywatny kanał przekazywania
transakcji do górników~\cite{flashbots}, część konkurencji przeniosła się poza publiczną pulę
transakcji, co dodatkowo utrudnia ocenę wykonalności na podstawie samych danych on-chain.
Propozycje ograniczania takich manipulacji na poziomie protokołu omawiają m.in. Heimbach
i~Wattenhofer~\cite{heimbach_sandwich}.

Badania empiryczne potwierdzają, że rozbieżności między pulami występują i~są konsumowane na dużą
skalę. McLaughlin i~in.~\cite{mclaughlin_arbitrage} w~pomiarze obejmującym 28~miesięcy
zidentyfikowali 3{,}8~mln wykonanych arbitraży o~łącznym zysku ok.~321~mln USD. Zaobserwowali też,
że część okazji utrzymuje się przez kilka kolejnych bloków. Berg i~in.~\cite{berg_inefficiencies}
analizowali pary Uniswap i~Sushiswap i~wykazali, że efektywność cenowa obu giełd spada w~okresach
wysokiej zmienności kursów, a~liczba okazji arbitrażu cyklicznego wówczas rośnie. Metodę
identyfikacji arbitraży cyklicznych w~Uniswap~V2 z~danych on-chain opisali Wang
i~in.~\cite{wang_cyclic_arbitrage}. Z~tych obserwacji wynikają dwie decyzje projektowe: dobór okien
historycznych obejmujących gwałtowne ruchy cen oraz śledzenie konsumpcji okazji przez kilka
bloków po jej wykryciu (rozdział~3.7).

\subsection{Logika rozmyta i~systemy wnioskowania rozmytego}

Zbiór rozmyty $\tilde{A}$ w~przestrzeni $U$ jest określony funkcją przynależności
$\mu_{\tilde A}\colon U \to [0,1]$~\cite{zadeh1965}. W~odróżnieniu od zbioru klasycznego element
może należeć do niego częściowo. Pojęcia nieostre, takie jak ,,duża rozbieżność'' czy ,,zaporowy
koszt gazu'', da się wtedy opisać funkcjami ciągłymi zamiast sztywnych progów. Dla oceny
wykonalności jest to naturalne: poszczególne czynniki nie mają ostrych granic i~wzajemnie się
kompensują.

System wnioskowania rozmytego typu Mamdaniego~\cite{mamdani1975} składa się z~bazy reguł postaci
,,\textsc{jeżeli} $x_1$ jest $\tilde A_1$ \textsc{i} $x_2$ jest $\tilde A_2$ \textsc{to} $y$ jest
$\tilde B$'', bloku rozmywania wejść, mechanizmu wnioskowania i~bloku wyostrzania. Stopień
aktywacji reguły oblicza t-norma (w~pracy minimum) stopni przynależności przesłanek. Konkluzja
reguły jest obcinana na tym poziomie (implikacja minimum), konkluzje wszystkich reguł agreguje
s-norma (maksimum), a~wynik liczbowy wyznacza środek ciężkości zagregowanego zbioru:
\begin{equation}
  y^{*} = \frac{\int y\, \mu_{\mathrm{agg}}(y)\, \mathrm{d}y}{\int \mu_{\mathrm{agg}}(y)\, \mathrm{d}y}.
  \label{eq:centroid}
\end{equation}
Zaletą modelu Mamdaniego jest interpretowalność. Reguły zapisują wiedzę ekspercką wprost, można
je czytać i~poprawiać bez danych uczących. Wadą jest brak mechanizmu strojenia: parametry funkcji
przynależności i~zestaw reguł dobiera się ręcznie. System Mamdaniego jako doradcę transakcyjnego
na rynku kryptowalut, z~regułami strojonymi algorytmem genetycznym, przedstawili Tupinambás
i~in.~\cite{tupinambas_fuzzy_crypto}.

Model Takagiego--Sugeno~\cite{takagi_sugeno1985} zachowuje rozmyte przesłanki, ale konkluzją reguły
$i$ jest funkcja wejść, w~wariancie pierwszego rzędu liniowa:
\begin{equation}
  y_i = c_{i0} + \sum_{j} c_{ij}\, x_j ,
  \qquad
  y = \frac{\sum_i w_i\, y_i}{\sum_i w_i},
  \label{eq:ts}
\end{equation}
gdzie $w_i$ jest stopniem aktywacji reguły. Wyjście~\eqref{eq:ts} jest ważoną średnią konkluzji,
przez co model jest różniczkowalny względem parametrów. Jang~\cite{jang1993} zaproponował ANFIS
(ang. \emph{adaptive-network-based fuzzy inference system}): model Takagiego--Sugeno zapisany jako
pięciowarstwowa sieć adaptacyjna (rozmywanie, iloczyn przesłanek, normalizacja aktywacji,
konkluzje liniowe, sumowanie) i~uczony algorytmem hybrydowym. Parametry konkluzji $c_{ij}$
wyznacza metoda najmniejszych kwadratów przy ustalonych przesłankach, a~parametry funkcji
przynależności koryguje krok gradientowy. ANFIS łączy czytelną strukturę regułową ze zdolnością
dopasowania do danych, kosztem częściowej utraty interpretowalności konkluzji. W~pracy oba
podejścia porównano na tym samym zestawie cech: system Mamdaniego z~regułami zapisanymi na
podstawie wiedzy dziedzinowej oraz ANFIS uczony na etykietach z~weryfikacji retrospektywnej
(rozdział~3.6).

\subsection{Ocena klasyfikatorów przy silnej nierównowadze klas}

Modele oceny wykonalności zwracają wynik liczbowy, który po przyłożeniu progu staje się decyzją
binarną. Ocena takich modeli wymaga miar niezależnych od progu i~odpornych na nierównowagę klas.
W~analizowanych danych okazje faktycznie skonsumowane stanowią około jednej dziesiątej
wykrytych rozbieżności.

Krzywa ROC (ang. \emph{receiver operating characteristic}) przedstawia zależność czułości od odsetka
fałszywych alarmów dla wszystkich progów. Pole pod nią (AUC) ma interpretację prawdopodobieństwa,
że losowo wybrany przykład pozytywny otrzyma wyższy wynik niż losowo wybrany
negatywny~\cite{fawcett2006}. Przy skośnym rozkładzie klas ROC bywa zbyt optymistyczna, bo duża
liczba negatywów maskuje fałszywe alarmy. Davis i~Goadrich~\cite{davis_goadrich2006} oraz Saito
i~Rehmsmeier~\cite{saito_pr_auc} pokazują, że krzywa precyzja--czułość i~pole pod nią (PR-AUC,
równoważne średniej precyzji) lepiej oddają użyteczność modelu, gdy liczy się jakość wskazań
pozytywnych. W~rozdziale~4 raportowane są obie miary, uzupełnione o~precyzję, czułość i~miarę
$F_1$ przy stałym progu. Niepewność AUC wyznaczono metodą bootstrap~\cite{efron1979}:
wielokrotne losowanie ze zwracaniem zbioru testowego daje rozkład statystyki, z~którego
odczytuje się przedział ufności. Sąsiednie bloki tej samej okazji są silnie skorelowane, a~jedna
transakcja może konsumować kilka rozbieżności, dlatego podział na zbiory uczący i~testowy musi
przebiegać na poziomie grup, nie pojedynczych obserwacji. Protokół opisano w~rozdziale~4.2.

\subsection{Istniejące rozwiązania i~uzasadnienie własnego podejścia}

Prace najbliższe tematycznie dzielą się na trzy grupy. Pierwszą stanowią pomiary ex post
wykonanych arbitraży i~MEV na podstawie historii łańcucha. McLaughlin
i~in.~\cite{mclaughlin_arbitrage} zbudowali system identyfikujący zrealizowane arbitraże oraz
detektor okazji obsługujący złożone modele cen; Qin i~in.~\cite{qin_bev} skwantyfikowali wartość
wydobytą z~różnych klas transakcji; narzędzie \emph{mev-inspect} Flashbots~\cite{mev_inspect}
klasyfikuje transakcje MEV w~zadanym zakresie bloków. Drugą grupę tworzą badania efektywności
rynku: Berg i~in.~\cite{berg_inefficiencies} mierzą, jak dokładnie Uniswap i~Sushiswap śledzą cenę
referencyjną, a~Wang i~in.~\cite{wang_cyclic_arbitrage} badają arbitraż cykliczny. Trzecią są
prace teoretyczne o~optymalnym arbitrażu w~rynkach o~stałej funkcji~\cite{angeris_uniswap}.

Wszystkie te prace odpowiadają na pytanie, co się wydarzyło albo jaki zysk był teoretycznie
możliwy. Żadna nie formułuje oceny wykonalności okazji \emph{ex ante}, w~chwili jej wykrycia, na
podstawie cech obserwowalnych w~danym bloku, ani nie weryfikuje takiej oceny względem
faktycznej konsumpcji. Tę lukę wypełnia system z~rozdziału~3. Łączy on wykrywanie rozbieżności
z~retrospektywną weryfikacją, czy zostały skonsumowane, a~uzyskane etykiety wykorzystuje do
porównania czterech modeli oceny (prostej heurystyki, jej wersji skalibrowanej, systemu Mamdaniego
i~ANFIS) na wspólnym protokole ewaluacji. Wnioskowanie rozmyte wybrano na rdzeń oceny z~powodu
charakteru problemu: czynniki wykonalności (rozbieżność, koszt gazu, płynność, ryzyko
konkurencji) są nieostre i~wzajemnie się kompensują, a~interpretowalność reguł pozwala wyjaśnić
każdą ocenę użytkownikowi narzędzia.
